% Generated from completed artifacts by prepare_results.py.
\newcommand{\DiffusionVerdict}{Diffusion scores lower than both JEPA checkpoints on the shared final-token chemical readout.}
\newcommand{\DiffusionPrimaryNumber}{0.7244}
\newcommand{\DiffusionTakeaway}{A matched token-level evaluation; training exposure and native input quantization still differ.}
\newcommand{\ReadoutTitle}{Diffusion trails the final-token chemical readout}
\newcommand{\ReadoutTakeaway}{Diffusion attention gains +0.0016 MACCS AUROC over mean; its paired interval includes zero.}
\newcommand{\MaccsTargetCount}{145}
\newcommand{\FunctionalTargetCount}{13}
\newcommand{\BootstrapCount}{1,000}
\newcommand{\BootstrapDefined}{980}
\newcommand{\PairedTitle}{Paired molecules separate the checkpoint differences}
\newcommand{\PairedTakeaway}{Attention differences: parent -0.0744; periodic -0.0506 MACCS AUROC.}
\newcommand{\SaeTitle}{Sparse reconstruction and chemical recognition are separate}
\newcommand{\SaeTakeaway}{Diffusion k=32 seed0: 92.8\% held-out variance explained; 3 chemical candidates frozen before test.}
\newcommand{\ChemicalCandidateSummary}{The strongest hydroxyl candidate (f4682) has test AUROC 0.744, but also mass correlation 0.599; strict matched-pair ordering is 0.517.}
\newcommand{\RelativeLTwoNote}{original-unit relative L2 is 0.259 for diffusion versus 0.088 / 0.054 for parent / periodic.}
\newcommand{\ClosingTakeaway}{Token-level readouts answer chemical accessibility; SAE structure adds a separate, qualified view.}
